% ----------------------------------------------------
% Projeto: UnDF-TeX
% Arquivo: 01-introducao.tex
%
% Copyright (c) 2026
% Universidade do Distrito Federal (UnDF)
%
% Este arquivo faz parte do projeto UnDF-TeX.
%
% Licenciado sob a LaTeX Project Public License (LPPL),
% versão 1.3c ou posterior.
%
% Implementação técnica baseada nas normas:
% ABNT NBR 14724:2024
% ABNT NBR 10520:2023
%
% Este projeto não reproduz o texto integral
% das normas da ABNT.
% ----------------------------------------------------

\chapter{Introdução}

Conforme \textcite{undfsite2026}, a Universidade do Distrito Federal disponibiliza informações institucionais em ambiente digital. A legislação educacional brasileira estabelece bases gerais para a organização do ensino \parencite{brasil1996ldb}.

Neste trabalho, o próprio TCC é tomado como objeto de explicação. O objetivo é mostrar, em um único documento, como se organizam os elementos obrigatórios e opcionais, como se estruturam os capítulos e como se apresenta o texto acadêmico de maneira formal.

\section{Contextualização do tema}

O Trabalho de Conclusão de Curso reúne pesquisa, escrita, revisão e apresentação estruturada do conhecimento. Por isso, compreender a forma do TCC é parte importante da própria experiência acadêmica.

\section{Objetivos}

Este exemplo completo tem como objetivo funcionar como modelo visual e estrutural de TCC, ajudando o estudante a compreender a aparência final do trabalho e a função de cada parte que o compõe.